\section{Discussion}
\label{sec:discussion}
The practical choice begins with the intended use of the samples. A
teacher-relative comparison is appropriate when preserving a generator's
conditional behavior under acceleration; observed-data regions are needed
when the target is representativeness of a measured population. A regional
repair is then assessed at its declared tolerance, with separate checks
for changes elsewhere in the distribution. The semantic continuation
illustrates this distinction: target-digit equivalence is established while
all-class equivalence remains unresolved. Reporting both makes the achieved
repair specific without expanding it into an unsupported distribution-wide
claim.

The comparison unit also determines what an apparent gain means. A
higher-order solver spends its calls on fewer integration intervals; a
prompt policy spends additional calls before generating its eventual output;
and a training continuation uses updates that a sampling-only remedy avoids.
The reported studies expose those costs separately. Their natural role is
to support a choice among measured quality, coverage, and resource tradeoffs
for the intended repeated use of a configuration.

The evidence also defines the operating scope. Exact regional guarantees
apply to fixed events at the declared effect size; approximate spread and
training-seed intervals have separate assumptions. Validated semantic regions
and learned fidelity scores retain evaluator limitations. The semantic repair
has three training seeds; natural-image \mbox{continuation} has one adaptation
seed. Frozen policy comparisons expose diagnostic overhead and unresolved
matched-cost advantages. Extending the confirmed repairs and reducing the
sample cost of reliable regional diagnosis are concrete next steps.
